\documentclass[12pt]{article}
\usepackage[T1]{fontenc}
\usepackage[utf8]{inputenc}
\usepackage{lmodern}
\usepackage{textcomp}
\usepackage{enumitem}
\usepackage{geometry}
\usepackage{graphicx}
\usepackage{amsmath, amssymb}
\geometry{margin=1in}
\begin{document}
\begin{center}
History - Undergraduate Level Assessment 8 \
Total Marks: 40 \
Answer all questions.
\end{center}

\section*{Multiple Choice (10 marks)}
\begin{enumerate}[label=\arabic*.]
\item What was NOT a deciding factor in the development of Mesopotamian civilization?
\begin{enumerate}[label=(\alph*)]
    \item the need to form alliances with neighboring civilizations for trade and protection
    \item the need to develop a calendar system for agricultural planning
    \item the need to develop a complex social system that would allow the construction of canals
    \item the use of irrigation to produce enough food
    \item the creation of a military force for defense and conquest
\end{enumerate}
\item According to one theory, the power of Ancestral Pueblo elites emerged when certain groups:
\begin{enumerate}[label=(\alph*)]
    \item created intricate pottery and jewelry.
    \item discovered the use of fire.
    \item invented the wheel.
    \item stored food surpluses in good years for distribution in bad years.
    \item established trade routes with distant cultures.
\end{enumerate}
\item Subsistence practices of people can be determined by recovering and studying:
\begin{enumerate}[label=(\alph*)]
    \item geological layers and sediment types.
    \item the architectural style and layout of structures on the site.
    \item pumice and evidence of pyroclastic surges.
    \item the faunal assemblage from a site.
    \item an osteological comparative collection from a site.
\end{enumerate}
\item Which of the following describes a key change in hominids beginning at least as early as Homo erectus that is probably related to increasingly larger brain size?
\begin{enumerate}[label=(\alph*)]
    \item cranial reduction
    \item decreased dentition
    \item microcephaly
    \item supraorbital cortex
    \item encephalization
\end{enumerate}
\item The Shang Dynasty laid the foundation for a coherent development of Chinese civilization that lasted well into the:
\begin{enumerate}[label=(\alph*)]
    \item 17 th century.
    \item 14 th century.
    \item 18 th century.
    \item 21 st century.
    \item 5 th century.
\end{enumerate}
\end{enumerate}

\section*{True / False (10 marks)}
\textit{Answer True or False}
\begin{enumerate}[label=\arabic*.]
\setcounter{enumi}{5}
\item True or False: The climate has become warmer in Montana and continues to do so. The glaciers in Glacier National Park have receded and are predicted to melt away completely in a few decades. Many Montana cities set heat records during July 2007, the hottest month ever recorded in Montana.
\item True or False: OTS, previously located at Lackland AFB, Texas until 1993 and located at Maxwell Air Force Base in Montgomery, Alabama since 1993, in turn encompasses a different answer separate commissioning programs: Basic Officer Training (BOT), which is for line-officer candidates of the active-duty Air Force and the U.
\item True or False: The angled deck allows the installation of one or two "waist" catapults in addition to the two bow cats. a different answer also improves launch and recovery cycle flexibility with the option of simultaneous launching and recovery of aircraft.
\item True or False: Examples include the apparent motion of the sun across the sky, the phases of the moon, the swing of a pendulum, and the beat of a heart. Currently, the international unit of time, the second, is defined by measuring the electronic transition frequency of a different answer (see below).
\item True or False: a different answer of the Nicaraguan population are indigenous.
\end{enumerate}

\section*{Long Form Response (20 marks)}
\textit{Answer each question with a clear and complete explanation.}
\begin{enumerate}[label=\arabic*.]
\setcounter{enumi}{10}
\item Read the following passage and answer the question: The history of comics has followed different paths in different cultures. Scholars have posited a pre-history as far back as the Lascaux cave paintings. By the mid-20 th century, comics flourished particularly in the United States, western Europe (especially in France and Belgium), and Japan. The history of European comics is often traced to Rodolphe Töpffer's cartoon strips of the 1830 s, and became popular following the success in the 1930 s of strips and books such as The Adventures of Tintin. American comics emerged as a mass medium in the early 20 th century with the advent of newspaper comic strips; magazine-style comic books followed in the 1930 s, in which the superhero genre became prominent after Superman appeared in 1938. Histories of Japanese comics and cartooning (manga) propo... Question: Which superhero appeared in comics in 1938?
\item Read the following passage and answer the question: Several certification programs exist to support the professional aspirations of software testers and quality assurance specialists. No certification now offered actually requires the applicant to show their ability to test software. No certification is based on a widely accepted body of knowledge. This has led some to declare that the testing field is not ready for certification. Certification itself cannot measure an individual's productivity, their skill, or practical knowledge, and cannot guarantee their competence, or professionalism as a tester. Question: With several certifications out there that can be aquired, what is the one trait they all share?
\end{enumerate}

\vfill
\noindent\textit{End of Paper}
\end{document}